%%%%%%%%%%%%%%%%%%%%%%%%%%%%%%%%%%%%%%%%%
% Short Sectioned Assignment LaTeX Template Version 1.0 (5/5/12)
% This template has been downloaded from: http://www.LaTeXTemplates.com
% Original author:  Frits Wenneker (http://www.howtotex.com)
% License: CC BY-NC-SA 3.0 (http://creativecommons.org/licenses/by-nc-sa/3.0/)
%%%%%%%%%%%%%%%%%%%%%%%%%%%%%%%%%%%%%%%%%

% \documentclass[paper=a4, fontsize=11pt]{scrartcl} % A4 paper and 11pt font size
\documentclass[11pt, a4paper]{book}
\usepackage[T1]{fontenc} % Use 8-bit encoding that has 256 glyphs
\usepackage[utf8]{inputenc}
\usepackage{fourier} % Use the Adobe Utopia font for the document - comment this line to return to the LaTeX default
\usepackage{listings} % para insertar código con formato similar al editor
\usepackage[spanish, es-tabla]{babel} % Selecciona el español para palabras introducidas automáticamente, p.ej. "septiembre" en la fecha y especifica que se use la palabra Tabla en vez de Cuadro
\usepackage{url} % ,href} %para incluir URLs e hipervínculos dentro del texto (aunque hay que instalar href)
\usepackage{graphics,graphicx, float} %para incluir imágenes y colocarlas
\usepackage[gen]{eurosym} %para incluir el símbolo del euro
\usepackage{cite} %para incluir citas del archivo <nombre>.bib
\usepackage{enumerate}
\usepackage{hyperref}
\usepackage{graphicx}
\usepackage{tabularx}
\usepackage{booktabs}

\usepackage[table,xcdraw]{xcolor}
\hypersetup{
	colorlinks=true,	% false: boxed links; true: colored links
	linkcolor=black,	% color of internal links
	urlcolor=cyan		% color of external links
}
\renewcommand{\familydefault}{\sfdefault}
\usepackage{fancyhdr} % Custom headers and footers
\pagestyle{fancyplain} % Makes all pages in the document conform to the custom headers and footers
\fancyhead[L]{} % Empty left header
\fancyhead[C]{} % Empty center header
\fancyhead[R]{Lucía González Sánchez} % My name
\fancyfoot[L]{} % Empty left footer
\fancyfoot[C]{} % Empty center footer
\fancyfoot[R]{\thepage} % Page numbering for right footer
%\renewcommand{\headrulewidth}{0pt} % Remove header underlines
\renewcommand{\footrulewidth}{0pt} % Remove footer underlines
\setlength{\headheight}{13.6pt} % Customize the height of the header

\usepackage{titlesec, blindtext, color}
\definecolor{gray75}{gray}{0.75}
\newcommand{\hsp}{\hspace{20pt}}
\titleformat{\chapter}[hang]{\Huge\bfseries}{\thechapter\hsp\textcolor{gray75}{|}\hsp}{0pt}{\Huge\bfseries}
\setcounter{secnumdepth}{4}
\usepackage[Lenny]{fncychap}


\begin{document}

	% Plantilla portada UGR
	\begin{titlepage}
\newlength{\centeroffset}
\setlength{\centeroffset}{-0.5\oddsidemargin}
\addtolength{\centeroffset}{0.5\evensidemargin}
\thispagestyle{empty}

\noindent\hspace*{\centeroffset}\begin{minipage}{\textwidth}

\centering
\includegraphics[width=0.9\textwidth]{logos/logo_ugr.jpg}\\[1.4cm]

\textsc{ \Large TRABAJO FIN DE GRADO\\[0.2cm]}
\textsc{ GRADO EN INGENIERIA INFORMATICA}\\[1cm]

{\Huge\bfseries Heráldica para todas \\}
\noindent\rule[-1ex]{\textwidth}{3pt}\\[3.5ex]
{\large\bfseries Subtítulo }
\end{minipage}

\vspace{2.5cm}
\noindent\hspace*{\centeroffset}
\begin{minipage}{\textwidth}
\centering

\textbf{Autor}\\ {Lucía González Sánchez}\\[2.5ex]
\textbf{Director}\\ {Juan Julián Merelo Guervós}\\[2cm]
\includegraphics[width=0.3\textwidth]{logos/etsiit_logo.png}\\[0.1cm]
\textsc{Escuela Técnica Superior de Ingenierías Informática y de Telecomunicación}\\
\textsc{---}\\
Granada, \today
\end{minipage}
\end{titlepage}


	% Plantilla prefacio UGR
	\thispagestyle{empty}

\begin{center}
{\large\bfseries Heráldica para todas \\ Subtítulo }\\
\end{center}
\begin{center}
Lucía González Sánchez\\
\end{center}

%\vspace{0.7cm}

\vspace{0.5cm}
\noindent\textbf{Palabras clave}: \textit{software libre}
\vspace{0.7cm}

\noindent\textbf{Resumen}\\
	

\cleardoublepage

\begin{center}
	Lucía González Sánchez\\
\end{center}
\vspace{0.5cm}
\noindent\textbf{Keywords}: \textit{open source}, \textit{floss}
\vspace{0.7cm}

\noindent\textbf{Abstract}\\


\cleardoublepage

\thispagestyle{empty}

\noindent\rule[-1ex]{\textwidth}{2pt}\\[4.5ex]

D. \textbf{Tutora/e(s)}, Profesor(a) del ...

\vspace{0.5cm}

\textbf{Informo:}

\vspace{0.5cm}

Que el presente trabajo, titulado \textit{\textbf{Chief}},
ha sido realizado bajo mi supervisión por \textbf{Estudiante}, y autorizo la defensa de dicho trabajo ante el tribunal
que corresponda.

\vspace{0.5cm}

Y para que conste, expiden y firman el presente informe en Granada a Junio de 2018.

\vspace{1cm}

\textbf{El/la director(a)/es: }

\vspace{5cm}

\noindent \textbf{(nombre completo tutor/a/es)}

\chapter*{Agradecimientos}






	% Índice de contenidos
	\newpage
	\tableofcontents

	% Índice de imágenes y tablas
	\newpage
	\listoffigures

	% Si hay suficientes se incluirá dicho índice
	\listoftables 
	\newpage

	% Introducción 
	\chapter{Introducción}

En este capítulo se presenta la motivación del proyecto, así como la definición
del problema que busca abordar. Se incluyen además las historias de usuario que
justifican el desarrollo de la solución propuesta, los objetivos iniciales y la
estructura del documento.

\section{Motivación}
La heráldica es un sistema de representación gráfica codificado que permite
la representación visual de linajes, instituciones y territorios a través de
un conjunto no arbitrario de reglas de diseño y composición. Estas normas
establecidas la hacen un campo apropiado para el análisis estructurado y la
automatización mediante herramientas digitales.

Sin embargo, quienes desean introducirse en esta disciplina se enfrentan a
ciertos obstáculos que dificultan su acceso y comprensión a, debido a que su
documentación se encuentra dispersa y no está completamente digitalizada. Estas 
barreras complican su aprendizaje, aplicación académica y divulgación.

Con la finalidad de satisfacer las necesidades de este grupo de personas y hacer
de esa manera menos hostil el aprendizaje de esta ciencia-arte, este proyecto busca 
desarrollar una herramienta que modernice y agilice el aprendizaje de la heráldica, 
facilitando su comprensión y análisis, aspirando a ofrecer un recurso digital accesible
que contribuya a la preservación y divulgación de esta disciplina.

\section{Definición del problema}

La falta de una herramienta tecnológica que permita analizar y clasificar escudos de 
armas de manera eficiente impide que las personas que estudian o investigan la heráldica
puedan consultar, clasificar o comparar escudos de forma eficiente, ya que en la actualidad
se requieren conocimientos especializados y un acceso a fuentes dispersas, muchas de ellas
no digitalizadas.

Esta situación presenta varias dificultades:

\begin{itemize}
    \item La consulta de escudos requiere revisar múltiples fuentes no centralizadas, 
    en su mayoría analógicas.
    \item No existen bases de datos estandarizadas que permitan buscar escudos por
    características visuales o estructurales.
    \item El análisis se realiza manualmente, lo que ralentiza el proceso y puede
    generar errores.
\end{itemize}

El objetivo de este proyecto será dar solución a los problemas presentados por los
clientes en las siguientes historias de usuario:

\subsection{Historias de usuario}
Con el fin de exponer y comprender mejor las necesidades a las que responde este
proyecto, se presentan los siguientes casos de uso:

\begin{itemize}
    \item \textbf{HU1: Estudiante de Historia del Arte - }Un estudiante de Historia del 
    Arte está paseando por Roma y se encuentra con una obra arquitectónica en cuya fachada
    hay un escudo heráldico. Se pregunta de quién será y qué representa. Sigue caminando 
    y se encuentra con más escudos, lo que le hace pensar en lo oportuna que sería una 
    plataforma que le ayude a identificarlos y obtener información sobre ellos de manera 
    sencilla.
    \item \textbf{HU2: Investigador en heráldica - }Un investigador encuentra tedioso y
    complicado analizar y comprobar escudos en fuentes analógicas, ya que hacerlo manualmente
    de esta manera lo convierte en un proceso lento y propenso a errores.
\end{itemize}

\section{Estructura del proyecto}
Es resto del documento se organiza en los siguientes capítulos:

\begin{itemize}
    \item \textbf{Capítulo 2: Estado del arte - }Se analizan trabajos previos relacionados
    con la digitalización de la heráldica y las herramientas existentes, estableciendo
    el contexto teórico del proyecto.
    \item \textbf{Capítulo 3: Planificación - }Se explica la planificación y metodología
    que se emplean en el desarrollo. Se habla de los milestones del proyecto.
    \item \textbf{Capítulo 4: Análisis - }Define qué debe hacer la herramienta y habla
    sobre su diseño y arquitectura.
    \item \textbf{Capítulo 5: Implementación - }Describe cómo se ha desarrollado y validado
    la herramienta. 
    \item \textbf{Capítulo 6: Conclusiones - }Se resumen los resultados obtenidos y se
    discuten posibles mejoras y líneas futuras de investigación.
\end{itemize}

Este proyecto es software libre, y está liberado con la licencia \cite{gplv3}.

	% Descripción del problema y hasta donde se llega
	\chapter{Descripción del problema}



	% Estado del arte
	% 	1. Crítica al estado del arte
	% 	2. Propuesta
	\chapter{Estado del arte}

El software libre y sus licencias \cite{gplv3} ha permitido llevar a cabo una expansión del
aprendizaje de la informática sin precedentes.

	
	\chapter{Planificación}

\section{Metodología utilizada}

\subsection{Herramientas utilizadas}

\subsubsection{Memoria}
Para la realización de esta memoria, se ha establecido como requisito mínimo que
el documento esté libre de faltas de ortografía y compile correctamente. Partiendo de que 
está escrito en \LaTeX, se han buscado herramientas que permitan la automatización de estas 
comprobaciones. Para ello se han estudiado dos tipos de herramientas: las destinadas a la
corrección gramatical y de estilo, y las empleadas para compilar la memoria.
\\

Entre las \textbf{herramientas para la corrección ortográfica}, han sido analizadas las siguientes
opciones:

\begin{itemize}
    \item \textbf{Textidote}: es una solución de código abierto diseñada específicamente para
    \LaTeX\, lo que evita falsos positivos por malinterpretar los comandos. Requiere Java.
    \item \textbf{LanguageTool}: esta opción de código abierto es más avanzada que la anterior 
    debido a que usa redes neuronales para comprobar la gramática y el estilo. 
    Sin embargo, no está diseñada específicamente para \LaTeX, lo que hace necesario 
    invertir tiempo extra en eliminar los comandos antes de realizar las verificaciones.
    Tiene una versión gratuita y una versión prémium con funcionalidades avanzadas.
    \item \textbf{GNU Aspell}: es el corrector ortográfico estándar de GNU, aunque compila para
    otros sistemas operativos. Es de código abierto, no corrige gramática ni estilo, solo
    ortografía, y no está diseñado específicamente para \LaTeX.
\end{itemize}

Después de la comparativa, la opción elegida ha sido \textbf{Textidote}, porque a pesar de 
no ser tan avanzada en la corrección, supone un ahorro considerable de tiempo el hecho de no 
tener que eliminar los comandos de \LaTeX\ al verificar la gramática y el estilo, como sucede
con las otras dos alternativas analizadas.
\\

Con respecto a las \textbf{herramientas para la compilación de la memoria}, se han estudiado:

\begin{itemize}
    \item \textbf{TeX Live}: es la distribución por defecto de \LaTeX\ para Linux. Sin embargo,
    la instalación puede ser pesada y excesiva si únicamente se necesita compilar documentos
    sencillos.
    \item \textbf{Overleaf}: es un editor \LaTeX\ basado en la nube, por lo que no requiere 
    instalación previa. Es ideal para trabajar de manera colaborativa al permitir
    a varios usuarios trabajar en el mismo documento y ver los cambios a tiempo real.
    No obstante, tiene limitaciones en cuanto a capacidad, que pueden solventarse con
    la versión de pago.
    \item \textbf{TinyTex}: es una versión basada en Tex Live, pero mucho más ligera. Está orientada
    sobre todo a usuarios de R, aunque es posible instalar manualmente los paquetes 
    específicos que sean necesarios, lo que la convierte en una buena opción en un contexto 
    en el que los recursos están limitados.
\end{itemize}

Dado a la complejidad añadida de instalar manualmente los paquetes específicos necesarios
para el proyecto y que los recursos no son limitantes en este caso, la herramienta elegida
es \textbf{TeX Live}. Se descarta \textbf{Overleaf} debido a la limitación de la versión gratuita,
además de que no será necesario un entorno de trabajo colaborativo.

\section{Temporización}

\section{Seguimiento del desarrollo}


	% Análisis del problema
	% 1. Análisis de requisitos
	% 2. Análisis de las soluciones
	% 3. Solucion propuesta
	% 4. Análisis de seguridad
	\chapter{Análisis del problema}
 


	% Desarrollo bajo sprints: 
	% 	1. Permitir registros y login de usuarios
	% 	2. Desarrollo del sistema de incidencias
	% 	3. Desarrollo del sistema de denuncias administrativas y accidentes
	% 	4. Desarrollo del sistema de croquis
	%   5. Instalación de la aplicación de manera automática
	\chapter{Implementación}

La implementación del software se ha dividido en hitos. Estos, han sido definidos en Github
y cada uno de ellos contiene un grupo de \textit{issues} que se corresponden con las distintas
mejoras que se han ido incorporando al software a lo largo de su desarrollo.\\



	% Presupuesto

	% Conclusiones
	\chapter{Conclusiones y trabajos futuros}



	% Trabajos futuros


	
	\newpage
	\bibliography{bibliografia}
	\bibliographystyle{plain}
	
\end{document}

